%=====================================================================================
\section{Conclusion}\label{sec:conclusion}

Can an applied researcher trust an LLM stance score of FOMC statements, and can they tell whether the
model read the text or remembered it? For the frontier model studied here, the answer to the first
question is largely yes, as a matter of measurement: its scores barely depend on the wording of the
prompt, and they survive the removal of the rate decision. That consistency may itself owe something to
recognition, and the few statements the model probably has not seen are too few to tell. The answer to
the second question is no, at least with the tools tested here. The model identifies almost every
statement even after blinding and paraphrasing, and it sometimes cites later events while scoring. Its
scores are weakly related to monetary policy surprises at no-change meetings, but that relation does not
survive a multiple-testing correction, is concentrated in recent years, and cannot be attributed to
reading rather than remembering.

Four practical lessons follow for work that uses LLMs to measure text.
\begin{enumerate}
\item Report reliability where it matters. Agreement on raw levels can be driven by the obvious content
(here, the rate decision). Report agreement after removing it and on changes, and report variation across
models as well as prompts: for changes, the model mattered about twice as much as the prompt.
\item Test recognition directly. A dating probe calibrated against a text-similarity baseline, and a
search of the scoring replies for unprompted identification, are cheap and should accompany any LLM
measure of well-known documents.
\item Do not rely on anonymisation or paraphrase to rule out look-ahead bias for well-known documents.
They do not prevent identification. Only models trained on data that stop before the documents, or
documents the model cannot have seen, can rule it out.
\item Treat a frontier model's market validation on well-known documents as possibly contaminated until such
a test confirms it, and report the specification search, the placebo distribution of the test statistic
and the influence of individual observations, calibrated against what a real effect would show.
\end{enumerate}

\paragraph{Limitations.} The study uses one frontier model; another may behave differently, and the
dating probe was not run on a second frontier model. The crossed
reliability block covers \R{icc.fable.n} statements from 2015--2019. Replies were stored cut at 300
characters until October 2026, so the counts of unprompted recognition are lower bounds and some long
replies could not be re-parsed. The two latest statements could not be checked against the Federal
Reserve's website. The paper uses no fine-tuned text classifier (such as FinBERT) as a further baseline,
does not split the surprises into the forward-guidance and asset-purchase factors of
\citet{swanson2021}, and has no human labels of meeting-to-meeting changes. It does not control for the Chair's press
conferences, which affect the daily yield changes, and it does not compare the scores with the text
measures of \citet{lucca2009}, \citet{hansen2016} or \citet{gardner2022}. Finally, the models were
accessed through a commercial interface in 2026; providers can change the model
behind a name, so exact replication of the scores may not be possible later, although every number in
this paper can be recomputed from the stored scores.
